% Root file used ashow macros
\documentclass{beamer}
\usepackage{tikz}
\usetikzlibrary{arrows.meta}
\title{Year 10 Vector Starters}
\date{}

\usepackage{multicol}

% ================================================================
% Answer-overlay helpers
% ================================================================
\newcommand{\ablank}[1]{%
  \alt<2>{\textcolor{red}{#1}}{\underline{\phantom{#1}}}%
}
\newcommand{\ashow}[1]{\uncover<2->{\textcolor{red}{\small #1}}}
\newcommand{\ashowq}[1]{\alt<2->{\textcolor{red}{\small #1}}{?}}

% Answers drawn straight on to a TikZ diagram, revealed on overlay 2 like \ashow.
% Space is reserved on overlay 1, so the diagram does not move between slides.
\usetikzlibrary{overlay-beamer-styles}
\tikzset{
  ans/.style     = {red, thick, visible on=<2->},
  ansfill/.style = {red!15, visible on=<2->},
  anslab/.style  = {red, font=\tiny, visible on=<2->},
}

\begin{document}

\frame{\titlepage}

%------------------------------------------------
\begin{frame}{Contents}
\small
\begin{columns}[T]
\column{0.5\textwidth}
\textbf{Questions and short answers}
\begin{itemize}
  \item \hyperlink{q-st1}{\beamergotobutton{Starter 1}}
\end{itemize}

\column{0.5\textwidth}
\textbf{Worked solutions}
\begin{itemize}
  \item \hyperlink{work-st1}{\beamergotobutton{Starter 1}}
\end{itemize}
\end{columns}
\end{frame}

%------------------------------------------------
\begin{frame}{Starter 1}
\hypertarget{q-st1}{}

\begin{multicols}{2}
\begin{enumerate}
  \item Start with zero, add 7, subtract 10, add 6, subtract 3, what is the result? \ashow{0}
  \item What if you started with $-5$? \ashow{$-5$}
  \item Add:
  $\begin{pmatrix}4\\1\end{pmatrix}
  +
  \begin{pmatrix}2\\3\end{pmatrix}$ \ashow{\(\binom{6}{4}\)}
  \item Find:
  $-\begin{pmatrix}5\\2\end{pmatrix}$ \ashow{\(\binom{-5}{-2}\)}
  \item What column vector represents the reverse of the following vector:
  $\begin{pmatrix}3\\-4\end{pmatrix}$? \ashow{\(\binom{-3}{4}\)}
\end{enumerate}
\end{multicols}

\vspace{0.3cm}



\begin{columns}[T]

\column{0.45\textwidth}
\begin{center}
\begin{tikzpicture}[scale=0.7]

\coordinate (A) at (0,0);
\coordinate (C) at (2,1);
\coordinate (B) at (4,2);

\draw[->] (A)--(C) node[midway,above] {$\mathbf{a}$};
\draw[->] (C)--(B) node[midway,above] {$\mathbf{b}$};

\draw[->,dashed] (A)--(B);

\node[below left] at (A) {A};
\node[above right] at (B) {B};

\end{tikzpicture}
\end{center}

\column{0.55\textwidth}

\begin{enumerate}
    \setcounter{enumi}{5}
  \item Write an expression in terms of $\bf{a}$ and $\bf{b}$ that represents the whole journey from A to B. \ashow{\(\mathbf{a}+\mathbf{b}\)}
  \item Which vector would undo that journey? How many different ways can you write it? \ashow{\(\overrightarrow{BA} = -(\mathbf{a}+\mathbf{b}) = -\mathbf{a}-\mathbf{b}\)}
\end{enumerate}

\end{columns}

\end{frame}

%------------------------------------------------
\begin{frame}{Worked Solution: Starter 1, Question 1}
\hypertarget{work-st1}{}
\textbf{Question:} Start with zero, add 7, subtract 10, add 6, subtract 3. What is the result?

\textbf{Worked Solution:}
\[
0+7-10+6-3 = 7-10+6-3 = -3+6-3 = 3-3 = 0.
\]
So the result is \(0\).

\end{frame}

%------------------------------------------------
\begin{frame}{Worked Solution: Starter 1, Question 2}
\textbf{Question:} What if you started with $-5$?

\textbf{Worked Solution:}
\[
-5+7-10+6-3 = 2-10+6-3 = -8+6-3 = -2-3 = -5.
\]
The operations have total change \(0\), so the starting value is unchanged; the result is \(-5\).

\end{frame}

%------------------------------------------------
\begin{frame}{Worked Solution: Starter 1, Question 3}
\textbf{Question:} Add \(\begin{pmatrix}4\\1\end{pmatrix} + \begin{pmatrix}2\\3\end{pmatrix}\).

\textbf{Worked Solution:}
Add corresponding components:
\[
\begin{pmatrix}4\\1\end{pmatrix}
+
\begin{pmatrix}2\\3\end{pmatrix}
=
\begin{pmatrix}4+2\\1+3\end{pmatrix}
=
\begin{pmatrix}6\\4\end{pmatrix}.
\]

\end{frame}

%------------------------------------------------
\begin{frame}{Worked Solution: Starter 1, Question 4}
\textbf{Question:} Find \(-\begin{pmatrix}5\\2\end{pmatrix}\).

\textbf{Worked Solution:}
Multiplying a vector by \(-1\) changes the sign of each component:
\[
-\begin{pmatrix}5\\2\end{pmatrix}
=
\begin{pmatrix}-5\\-2\end{pmatrix}.
\]

\end{frame}

%------------------------------------------------
\begin{frame}{Worked Solution: Starter 1, Question 5}
\textbf{Question:} What column vector represents the reverse of the vector \(\begin{pmatrix}3\\-4\end{pmatrix}\)?

\textbf{Worked Solution:}
The reverse vector has the same length but the opposite direction, so take the negative:
\[
-\begin{pmatrix}3\\-4\end{pmatrix}
=
\begin{pmatrix}-3\\4\end{pmatrix}.
\]
So the reverse vector is \(\begin{pmatrix}-3\\4\end{pmatrix}\).

\end{frame}

%------------------------------------------------
\begin{frame}{Worked Solution: Starter 1, Question 6}
\textbf{Question:} Write an expression in terms of \(\mathbf{a}\) and \(\mathbf{b}\) that represents the whole journey from A to B.

\textbf{Worked Solution:}
The diagram shows \(\overrightarrow{AC}=\mathbf{a}\) and \(\overrightarrow{CB}=\mathbf{b}\). The journey from A to B goes via C, so the vectors are added tip-to-tail:
\[
\overrightarrow{AB}
=
\overrightarrow{AC}+\overrightarrow{CB}
=
\mathbf{a}+\mathbf{b}.
\]
So the required expression is \(\mathbf{a}+\mathbf{b}\).

\end{frame}

%------------------------------------------------
\begin{frame}{Worked Solution: Starter 1, Question 7}
\textbf{Question:} Which vector would undo that journey? How many different ways can you write it?

\textbf{Worked Solution:}
The vector that undoes \(\overrightarrow{AB}\) is its reverse, \(\overrightarrow{BA}\). Since \(\overrightarrow{AB}=\mathbf{a}+\mathbf{b}\),
\[
\overrightarrow{BA}
=
-\overrightarrow{AB}
=
-(\mathbf{a}+\mathbf{b})
=
-\mathbf{a}-\mathbf{b}.
\]
So the undo vector can be written in two equivalent forms:
\[
\overrightarrow{BA}
=
-(\mathbf{a}+\mathbf{b})
=
-\mathbf{a}-\mathbf{b}.
\]

\end{frame}

\end{document}